\section{Introduction}\label{introduction}

Consider a coding task from SWE-bench: resolve a Django migration bug.
The first phase is planning. An agent reads the issue, identifies the
affected modules, and decomposes the fix into three independent file
edits. Hierarchical orchestration handles this well; a planner delegates
subtasks to specialists, each working in parallel. But then the agent
runs the test suite and two tests fail for unexpected reasons. The task
has shifted. What started as a cleanly decomposable problem has become
an iterative debugging loop where multiple hypotheses need to be tested,
compared, and refined. Hierarchical orchestration keeps dispatching
subtasks to specialists who can't see each other's outputs. Debate,
where multiple agents argue about the root cause and converge through
critique, would be the better fit. But the system committed to hierarchy
at the start and has no way to switch.

This scenario plays out constantly in practice. Our prior work on
OrchestraBench \citep{anonymous_2026} showed that orchestration topology
choice is invisible on easy tasks but decisive on hard ones, with debate
reaching 58.6\% on hard tasks where flat conversation scored 18.6\%. The
decomposability-iterativeness-tool-diversity (D-I-T) framework we
developed predicts which topology fits which task with 90\% accuracy.
But that prediction happens once, before execution begins. It assumes
each task has a single dominant structure.

Real tasks don't hold still. A SWE-bench task might start with high
decomposability (plan the fix), shift to high iterativeness (debug the
failures), then return to moderate decomposability (write clean tests).
GAIA research tasks often begin with broad information gathering (high
tool diversity, low iterativeness) before narrowing into deep reasoning
(low tool diversity, high iterativeness). The structural characteristics
that determine optimal topology change as the task unfolds.

Every existing adaptive orchestration system misses this. Yu's AdaptOrch
framework \citep{adaptorch_2026} selects among four canonical topologies
at dispatch time based on task metadata. MDAgents \citep{kim_2024}
adapts between solo and group configurations for medical tasks, again
before execution starts. Evolving Orchestration \citep{evolving_2025}
trains an RL policy to sequence agents, but that policy is fixed once
trained. MASS \citep{mass_2025} and G-Designer \citep{gdesigner_2025}
optimize topologies offline. DyTopo \citep{dytopo_2026} reconstructs
communication graphs per reasoning round but does not switch between
structurally different topologies.

Meanwhile, a separate literature on execution trace analysis has shown
that rich structural signals exist in agent execution logs. MAST
\citep{cemri_2025} identified 14 failure modes from 150+ traces across 7
frameworks. TraceCoder \citep{tracecoder_2026} drives debugging cycles
from trace features. DIG \citep{dig_2026} captures collaboration
dynamics as time-evolving causal networks. These signals are used for
post-hoc diagnosis. Nobody feeds them back into topology decisions at
runtime.

We bridge these two literatures. Our system, which we call AdaptOrch-RT
(Adaptive Orchestration with Runtime Topology switching), monitors
execution traces in real time, estimates the current D-I-T
characteristics of the task, detects when those characteristics shift
past a threshold (a phase boundary), and switches to the topology whose
structural prior best matches the new phase. The switching decision is
formulated as a contextual bandit problem where the context is the
estimated D-I-T vector and the arms are three canonical topologies: flat
conversation, hierarchical decomposition, and multi-agent debate.

We make three contributions:

\begin{enumerate}
\item A \textbf{runtime D-I-T estimator} that extracts decomposability, iterativeness, and tool-diversity features from execution traces as they unfold. The estimator uses rolling windows over tool call diversity, subtask completion rate, and revision count to produce real-time structural characterizations.

\item A \textbf{switching policy} built on a contextual bandit that triggers topology transitions at detected phase boundaries. The policy includes a cost model that accounts for switching overhead (context transfer, re-prompting) and fires only when the expected benefit exceeds the transition cost.

\item \textbf{Controlled experiments} on the 82-task OrchestraBench suite showing that AdaptOrch-RT achieves 78\% task success versus 64.6\% for the best static topology (debate) and 90\% for the oracle that always picks the single best topology per task. Switching occurs on average 1.3 times per hard task, with 85\% of switches improving outcomes.
\end{enumerate}

The gap between 78\% and 90\% tells us where the remaining challenge
lies: predicting phase boundaries accurately and reducing switching
overhead. But the gap between 78\% and 64.6\% tells us that runtime
topology switching already recovers a substantial fraction of the
performance left on the table by static assignment.
